\section{Des collections d'exception au service du droit}

La bibliothèque Cujas regorge d'une richesse documentaire géré par le Département du Patrimoine, de la conservation et de la numérisation (DPCN) qui gère \enquote{40 0000 documents patrimoniaux du XVe au XXe siècle conservés en salle du Patrimoine}. Elle est forte d'un des plus importants fonds juridiques à l'échelle internationale.

La richesse des collections de Cujas repose sur des fonds spécialisés et des ouvrages rares. Cujas est connu pour conserver un \enquote{ensemble unique de ressources, de collections en sciences juridiques, économiques et politique}. Elle conserve aussi \enquote{les thèses, les mémoires et les publications officielles de la Société des Nations, des Nations Unies, du Journal officiel France et des anciennes colonnies}, ainsi qu'une \enquote{collection de travaux académiques imprimés en Allemagne aux XVIe et XIXe siècles(le fonds Jean-Marie Pardessus)}. Ces travaux retraçent la sociabilité universitaire et l'histoire de l'enseignement de l'ancienne Faculté de droit. La bibliothèque Cujas s'est efforcé de conserver \enquote{une collection quasi exhaustive de thèses de droit soutenues à la Faculté de droit}. Aujourd'hui encore, elle poursuit encore la collecte des thèses de droit, d'économie et de science politique soutenue à l'Université Paris I, Panthéon-Sorbonne et à l'Université Paris II, Panthéon-Assas. Le support de ces travaux varie selon la période passant de la forme imprimée, à la forme microfichée ou numérique.

Les collections de la bibliothèque Cujas ne s'arrêtent pas là et comprennent d'autres ressources comme les dissertations d'agrégation de droit, les manuels de pédagogie du droit. Les manuels de pédagogie du droit ont été constitué dans le cadre du projet HOPPE-Droit (Hypothèses d'Observation des Productions Pédagogiques Editées en Droit), financé par CollEx-Persée. Il s'agit d'un outil numérique permettant d'étudier l'évolution de l'enseignement du droit français du XIXe au XXe siècle. Grâce à la richesse de ces collections, Cujas profite de ce projet pour \enquote{sonder l'état de ses collections et enrichir ses données, tout en offrant aux chercheurs un instrument de travail inédit}. La particularité de l'outil HOPPE-Droit est de pouvoir croiser \enquote{des éléments des notices bibliographiques de chacun des ouvrages retenus dans le corpus qu'elle fait résonner les uns par rapport aux autres (titres, auteurs, domaine, éditeurs, collections)}. Cette approche permet de comprendre les évolutions du droit français et de son enseignement. L'outil permet d'utiliser ces données massives afin de voir les \enquote{tendances historiques qui se sont faites jour en matières d'ouvrages juridiques à vocation pédagogique}. Cet outil permet ainsi de mieux appréhender \enquote{les réseaux d'édition qui structurent le champ du droit et d'en suivre l'évolution dans le temps}. Ce projet a vu le jour grâce à la collaboration étroite des enseignants-chercheurs et des bibliothécaires avec comme objet d'étude celui de \enquote{suivre les évolutions de l'enseignement du droit, saisies à travers le prisme de l'édition juridique à vocation pédagogique}. 
	
Les fonds spécialisés de la bibliothèque Cujas ont été répartis en six catégories comprenant : les publications officielles de la Société des Nations (SDN), des Nations Unies (ONU), du Journal Officiel (JO), du fonds Jean-Marie Pardessus, des cartes postales et des archives de la Faculté de droit de Paris. La collections des publications officielles de la SDN a été constitué dans les années 1920, couvrant l'ensemble de ces publications entre 1919 et 1946. Cette collection comprend 534 volumes en français et en anglais. Pour les publications officiels de l'ONU, la bibliothèque Cujas n'est dépositaire du fonds que depuis 1949 constituant un fonds de près de 3 800 cartons. Toutefois, la bibliothèque a cessé le dépôt en 2010 avec la diffusion en ligne des documents officiels. Sur la période allant de 1870 à 2015, la bibliothèque conserve l'intégralité du JO de la République française. Les éditions papiers sont consultables sous format papier jusqu'en 2015 et c'est à partir de 2004 qu'il est possible de consulter leur version en ligne. Pour le fonds Pardessus, cette collection est constitué par Jean-Marie Pardessus (1772-1853) à travers laquelle il rassemble près de 9 000 travaux académique en droit parcours la période du XVIe jusqu'au XIXe siècle. Ils ont la particularité d'avoir \enquote{principalement été rédigés en latin et soutenues dans les universités allemandes}. Ce fonds entre en possession de la bibliothèque Cujas, grâce à la donation de M. de Rozière, le petit-fils de Jean-Marie Pardessus, à la Faculté de droit de Paris. La collection des cartes postales est un fond \enquote{caricaturales représentant des professeurs de la Faculté de droit de Paris au début du XXe siècle}. Elles constituent un fonds original pour l'histoire de l'enseignement du droit droit couvrant la période de 1900 à 1914. Le dernier fonds est celui des archives de la Faculté de droit de Paris. Ces archives retracent \enquote{la vie institutionnelle, pédagogique et politique de la Faculté de droit (XIXe-XXe siècle)}. Des sources administratives, des papiers de professeurs et des manuscrits composent ce fonds d'archives.
	
